% from ega4-17-fr.tex
\oldpage[IV-4]{59}
Si l'on tient compte de ce que les $i_\alpha$ sont formellement étales, tout revient donc à prouver que si les $f\circ i_\alpha$ sont formellement lisses (resp. formellement non ramifiés), il en est de même de $f$.

Soient donc $Y'$ un schéma affine, $Y'_0$ un sous-schéma fermé de $Y'$ défini par un Idéal nilpotent $\sh{J}$, que l'on peut supposer tel que $\sh{J}^2=0$ (17.1.2, (ii)), et enfin soit $g:Y'\to Y$ un morphisme. Supposons donné un $Y$-morphisme $u_0:Y'_0\to X$; désignons par $W_\alpha$ (resp. $W^0_\alpha$) le préschéma induit par $Y'$ (resp. $Y'_0$) sur l'ouvert $u_0^{-1}(U_\alpha)$ (on rappelle que $Y'$ et $Y'_0$ ont \emph{même espace topologique sous-jacent}). Supposons d'abord que les $f\circ i_\alpha$ soient \emph{formellement non ramifiés}, et montrons que, si $u'$ et $u''$ sont deux $Y$-morphismes de $Y'$ dans $X$ dont les restrictions à $Y'_0$ coïncident, alors on a $u'=u''$. En effet, compte tenu de (17.1.2, (iv)), l'hypothèse que les $f\circ i_\alpha$ sont non ramifiés entraîne que pour tout $\alpha$, on a $u'|W_\alpha=u''|W_\alpha$, puisque les restrictions de ces deux $Y$-morphismes à $W^0_\alpha$ coïncident. D'où la conclusion dans ce cas.

Supposons maintenant tous les $f\circ i_\alpha$ \emph{formellement lisses} et prouvons qu'il existe un $Y$-morphisme $u:Y'\to X$ dont $u_0$ est la restriction à $Y'_0$. Or, puisque $Y'$ est un \emph{schéma affine}, on peut appliquer (16.5.17) dont les hypothèses sont satisfaites, et dont la conclusion démontre précisément l'existence de $u$.

On peut donc dire que les notions introduites dans (17.1.1) sont \emph{locales} sur $X$ et sur $Y$, ce qui permet toujours, en vertu de (17.1.2, (i)), de se ramener à l'étude des \emph{algèbres} formellement lisses (resp. formellement non ramifiées, resp. formellement étales).

\subsection{Propriétés différentielles générales}
\EGAReaderTarget{ega4.c04.subsec.17-2}{subsection:IV.17.2}

\begin{proposition}[17.2.1]
\label{IV.17.2.1}
Pour qu'un morphisme $f:X\to Y$ soit formellement non ramifié, il faut et il suffit que $\Omega_f^1=0$ (ce qu'on écrit encore $\Omega_{X/Y}^1=0$ (16.3.1)).
\end{proposition}

Compte tenu de (17.1.6), on est ramené au cas où $Y=\Spec(A)$ et $X=\Spec(B)$ sont affines, et la conclusion résulte alors de (0, 20.7.4) et de l'interprétation de $\Omega_{X/Y}^1$ dans ce cas (16.3.7).

\begin{corollary}[17.2.2]
\label{IV.17.2.2}
Soient $f:X\to Y$, $g:Y\to Z$ deux morphismes. Pour que $f$ soit formellement non ramifié, il faut et il suffit que l'homomorphisme canonique (16.4.19)
\[
  f^*(\Omega_{Y/Z}^1)\longrightarrow\Omega_{X/Z}^1
\]
soit surjectif.
\end{corollary}

C'est une conséquence immédiate de (17.2.1) et de la suite exacte (16.4.19.1).

\begin{proposition}[17.2.3]
\label{IV.17.2.3}
Soit $f:X\to Y$ un morphisme formellement lisse.
\begin{enumerate}
  \item[(i)] Le $\sh{O}_X$-Module $\Omega_{X/Y}^1$ est localement projectif (16.10.1). Si $f$ est localement de type fini, $\Omega_{X/Y}^1$ est localement libre de type fini.
  \item[(ii)] Pour tout morphisme $g:Y\to Z$, la suite (16.4.19) de $\sh{O}_X$-Modules
  \[
  \label{IV.17.2.3.1}
    0\longrightarrow f^*(\Omega_{Y/Z}^1)\longrightarrow\Omega_{X/Z}^1\longrightarrow\Omega_{X/Y}^1\longrightarrow0
    \tag{17.2.3.1}
  \]
  est exacte; en outre, pour tout $x\in X$, il existe un voisinage ouvert $U$ de $x$ tel que les restrictions à $U$ des homomorphismes de (17.2.3.1) forment une suite exacte et scindée.
\end{enumerate}
\end{proposition}

